\section{Discussion}
\label{sec:discussion}

\paragraph{Cost and latency.}
VLM verification is the only paid stage. Escalation is bounded by its
trigger rate (51--58\% of objects), and per-call cost is $\sim$\$0.009;
full verification averages $\sim$\$0.09 per object (\$3.63 for a 38-cluster
run and \$7.37 for an 86-cluster run, all calls included). The incremental
design changes
the cost model qualitatively: a repeat scan costs \$0 (hash carry-over) and
a changed scene costs in proportion to its change (3.1\% of rebuild in our
measurement). Verification is offline; no VLM call sits on a robot's control
path.

\paragraph{Residual unverified content.}
On exploration-quality scans, 24\% of objects remain unverified (47\% on the
manual scan; 43\% on the 1/6-coverage ablation)---scan quality is the
dominant driver of verification yield. The framework's position is to expose
this residue, not shrink it cosmetically: unverified nodes stay queryable
with demoted types, and the gate is only as good as the verifier
(Sec.~5.3's unanimous-but-wrong cases; the escalated-verified keyboard trap
that survives gating on the robot hall). Verifier-bounded gating is a
property users must design around, and we report it as such.

\paragraph{Small low-reflectance objects.}
The clutter-absorption error class---disassembled robot, poster stand, and
the controlled-change floor circulator all verified as
\texttt{clutter}---shows a systematic blind spot: geometrically detected,
semantically flattened. Change detection still works (the circulator was
correctly inserted), but type-level recall for small dark appliances needs
either close-range imagery or targeted re-rendering.

\paragraph{Structure remnants.}
The soffit filter removes vertical transition bands, but horizontal and
sloped ceiling patches and content seen through openings survive structure
removal (six residual cases in the robot hall). These are bounded,
enumerable artifacts; extending the filter to non-vertical remnants is
mechanical future work.

\paragraph{The coverage paradox.}
Better scans can make density-based decomposition worse: at 99.6\% coverage,
gaps between adjacent objects close and DBSCAN chains 41\% of the scene into
one mega-cluster. Our measurements show this artifact accounts for only
$\sim$12\% of the epoch's absent nodes (5/43 recoverable by tighter
re-splitting; 4/43 by Point-SAM refinement)---the rest is genuine change---but
the effect argues for coverage-adaptive clustering or learned instance
priors as the backbone's next iteration, applied under the same
flagged-candidate discipline as our Point-SAM rule.

\paragraph{Scope.}
Results come from one scanner, one building, and object-level change;
structural change (walls, large fixtures) and cross-site generalization are
untested. The three-epoch study mitigates but does not remove the two-scene
limitation. The place layer's regions are parameterized SLIC superpixels over the grid
map: spatial cells rather than architecturally meaningful rooms, with the
superpixel count fixed by hand; inferring wall- and doorway-aligned region
boundaries is future work.
